%!TEX root = ../main.tex
% Resumen en español (~250 palabras). Las palabras clave se definen
% en main.tex con \palabrasclave{...}
%
% Guía de la cátedra (Entrega 0): el resumen responde, en este orden,
%   1. Problema: qué situación concreta de TI se aborda y a quién afecta.
%   2. Enfoque: qué se construye o investiga y con qué método o tecnología.
%   3. Valor: qué aporte original ofrece frente a lo existente.
% Impersonal, sin citas ni siglas sin expandir. Se reescribe al final del
% proyecto para que refleje lo que efectivamente se logró.

\begin{resumen}

GraphSAST es un sistema de análisis estático de código fuente (SAST)
orientado a la detección temprana de vulnerabilidades de arquitectura de
software. El sistema compila la lógica de aplicaciones escritas en
TypeScript/JavaScript en una representación formal mediante un grafo de
flujo de datos. A partir de técnicas de análisis de contaminación (taint
analysis), la herramienta rastrea el recorrido de datos no confiables
(sources) hacia operaciones sensibles (sinks) y verifica la presencia o
ausencia de mecanismos de sanitización en el trayecto. El proyecto abarca
el diseño e implementación propios del analizador, la construcción del
grafo, el motor de consultas y una interfaz web interactiva para la
visualización de los caminos de riesgo, validando su precisión y su
capacidad de reducción de falsos positivos frente a repositorios con
vulnerabilidades documentadas.

\end{resumen}
